% ============================================================================
%  abstract.tex — Résumé / Abstract Template
%  Replace the placeholder text below with your project's summary.
% ============================================================================

% --- French Abstract ---
\chapter*{Résumé}
\addcontentsline{toc}{chapter}{Résumé}

\begin{center}
{\large
Ce rapport présente les travaux réalisés dans le cadre du TP N°5 du module
\textbf{Programmation Orientée Objet en Java SE}, portant sur deux concepts
fondamentaux du langage Java~: la \textbf{gestion des exceptions} et la
\textbf{généricité}.\\[0.4cm]

Le TP est organisé en trois parties~: la première traite des exceptions
système (vérifiées et non vérifiées) à travers des exemples concrets de
division, de saisie utilisateur, de tableaux et de propagation d'exceptions.
La deuxième partie porte sur la création d'exceptions personnalisées
appliquées à des cas métier tels que les comptes bancaires, la gestion de
stock et les formulaires. La troisième partie explore la généricité avec
des classes, interfaces, méthodes génériques, bornes et wildcards.\\[0.4cm]

\textbf{Mots-clés~:} Java, Exceptions, Généricité, Checked, Unchecked,
Wildcards, Programmation Orientée Objet.
\vfill
}
\end{center}

% --- English Abstract (optional, uncomment if needed) ---
% \chapter*{Abstract}
% \addcontentsline{toc}{chapter}{Abstract}
%
% \begin{center}
% {\large
% Provide a concise summary of your project in English here.\\[0.4cm]
%
% State the problem, methodology, and key findings.\\[0.4cm]
%
% \textbf{Keywords:} Keyword~1, Keyword~2, Keyword~3, Keyword~4.
% \vfill
% }
% \end{center}
